\section{Задание}

Используя хвостовую рекурсию, разработать эффективную программу, позволяющую:
\begin{enumerate}
	\item Найти длину списка (по верхнему уровню);
	\item Найти сумму элементов числового списка;
	\item Найти сумму элементов числового списка, стоящих на нечетных позициях исходного
	      списка (нумерация от 0);
	\item Сформировать список из элементов числового списка, больших заданного значения;
	\item Удалить заданный элемент из списка (один или все вхождения).
	\item Объединить два списка.
\end{enumerate}

Убедиться в правильности результатов.
Для одного из вариантов вопроса уметь составить таблицу, отражающую конкретный
порядок работы системы.

\begin{tabularx}{\textwidth}{|p{1.1cm}|X|X|X|}
	\hline
	№ шага & Состояние резольвенты, и вывод: дальнейшие действия (почему?) & Для каких термов запускается алгоритм унификации: Т1=Т2 и каков результат (и подстановка) & Дальнейшие действия: прямой ход или откат (почему и к чему приводит?) \\
	\hline
	1....  & ...                                                           &                                                                                           & Комментарий, вывод...                                                 \\
	\hline
	2      & ...                                                           &                                                                                           & ...                                                                   \\
	\hline
\end{tabularx}

Т.к. резольвента хранится в виде стека, то состояние резольвенты требуется отображать в столбик: вершина -- сверху! Новый шаг надо начинать с нового состояния резольвенты!
